\section{서 론}
\label{sec:서론}

무인 비행체 애드혹 네트워크(Flying Ad-hoc Network; FANET)는 재난 구조, 군사 작전, 환경 모니터링 등 고정된 통신 인프라가 작동하지 않는 극단적 환경에서 신속하게 임시 통신망을 제공할 수 있는 유망한 기술로 평가받고 있다. 그러나 FANET을 구성하는 무인 항공기(Unmanned Aerial Vehicle; UAV)들은 공중에서 매우 빠른 속도로 불규칙하게 이동하기 때문에 무선 링크 상태가 동적으로 급격히 변화하며, 이로 인해 기존 무선 이동 애드혹 네트워크(MANET)보다 네트워크 토폴로지의 파괴와 형성이 훨씬 더 빈번하게 발생한다. 따라서 이러한 동적 FANET 환경 하에서 종단간 패킷 전달 성능(Packet Delivery Ratio; PDR)을 보장하고 전송 지연을 최소화하는 신뢰성 있는 분산 라우팅 프로토콜 설계는 핵심적인 당면 과제이다.

기존의 대표적인 FANET 라우팅 프로토콜로는 목적지 노드와의 물리적 좌표 거리가 가장 좁혀지는 방향의 이웃 노드로 패킷을 포워딩하는 GPSR(Greedy Perimeter Stateless Routing)\cite{gpsr}과 같은 지리적 라우팅(Geographic Routing) 기법이 널리 채택되어 왔다. 그러나 지리적 라우팅 기법은 장애물이 존재하거나 노드 배치가 불균일한 라우팅 홀(Routing Hole) 지역에 봉착했을 때, 로컬 미니멈(Local Minimum) 문제로 인해 경로 탐색에 실패하고 패킷을 최종 드롭하는 치명적인 취약점을 안고 있다. 이를 보완하기 위해 수치적 지리 진행도와 링크 속도 등을 활용한 예측적 지리적 라우팅 기법들이 고안되었으나, 휴리스틱 매개변수의 조율 한계와 무선 채널의 비정상성(Non-stationarity)으로 인해 복잡하게 우회해야 하는 환경에서는 일반화 성능을 충분히 보증하기 어렵다.

최근 이러한 제약을 기계학습 기반으로 해소하고자 다중 에이전트 강화학습(Multi-Agent Reinforcement Learning; MARL)을 결합한 라우팅 프로토콜이 적극적으로 제안되고 있다. 전역적인 네트워크 연결 상태를 모니터링하는 중앙 집중형 학습 및 분산 실행(Centralized Training and Decentralized Execution; CTDE) 기법들은 우수한 포워딩 행동 궤적을 스스로 학습할 수 있음이 확인되었다. 그러나 이러한 MARL 기법을 실제 제약된 자원의 UAV에 물리적으로 배포하는 단계에서는 심각한 \textbf{배포 격차(Deployment Gap)}가 발생한다. 구체적으로, 전역적인 그래프 주의 네트워크(GNN)나 잦은 제어 메시지 교환(Emergent Communication)에 의존하는 라우팅 기법들은 분산 환경의 각 UAV가 매 홉마다 엄청난 양의 무선 토폴로지 메시지를 실시간으로 교환하거나, 연산 비용이 높은 추론 연산을 지속적으로 수행해야 하므로 실용 배포가 거의 불가능하다. 즉, 실제 현장의 각 UAV는 자신의 1-hop 물리적 통신 영역 안에 진입한 후보 노드들의 지역적 관측 정보(Ego-graph)만을 기반으로 오버헤드 없이 신속한 분산 실행을 완료해야 한다.

본 논문에서는 이러한 전역 정보 기반 교사 정책의 우수한 지식을 로컬 관측만을 수행하는 가벼운 학생 정책으로 이전하고, 이상 징후 상황에서만 선택적으로 예측 기계학습 정책을 활성화하는 예측적 위험 스위칭 기반 분산형 라우팅 프레임워크인 \textbf{\method}를 제안한다. 제안 기법은 전역 정보를 완전하게 활용할 수 있는 가상의 교사(Teacher) GNN 정책을 오프라인에서 학습시킨 후, 그 행동 확률 분포(Action Probability Distribution) 자체를 로컬 1-hop 관측치만을 사용하는 다층 퍼셉트론(MLP) 구조의 학생(Student) 정책에 \textbf{정책 증류(Policy Distillation)}\cite{rusu2016policy} 기법을 사용하여 주입한다. 또한, 무선 페이딩 채널 손실 임계 및 버퍼 지연 위험 상태를 실시간 예측하는 스위칭 게이트(Predictive Risk Switch)를 도입하여, 평상시에는 초경량 지리 라우팅으로 오버헤드 없이 패킷을 흘려보내다가, 위험이 감지되는 예외 국면(라우팅 홀 진입, 예기치 못한 링크 단절)에서만 ML 기반 학생 정책의 우회 결정을 동적으로 구동한다.

본 논문의 핵심 기여(Contribution)는 다음과 같이 요약된다:
\begin{enumerate}
    \item \textbf{분산 의사결정 정식화}: FANET 내 각 UAV가 1-hop 이웃 상태만을 관측하여 차홉 결정을 내리는 실질적인 라우팅 문제를 부분관측 분산 의사결정 모형인 Dec-POMDP로 정교하게 정의하고 체계적인 제약을 정식화하였다.
    \item \textbf{정책 증류와 예측적 위험 스위칭 융합 프레임워크 제안}: 전역 GNN 교사의 포워딩 지식을 local MLP 학생에게 전수하는 정책 증류 기법과, 평시 지리 라우팅과 비상시 예측 라우팅 분기를 가볍게 상호 전환하여 연산/통신 비용을 최소화하는 선택적 위험 스위칭 구조(\method)를 제안하였다.
    \item \textbf{실제 데이터 기반 다각적 검증}: 총 14가지 동적 이동성 시나리오와 5가지 시드에 대한 84,000회 이상의 풀 시뮬레이션 검증을 통하여 제안 기법의 성능 우위(GPSR 대비 32.5\% PDR 향상, 고성능 MARL 대비 제어 오버헤드 바이트 대폭 절감)와 구조적 타당성을 엄밀하게 입증하였다.
\end{enumerate}
